% ------------------------------------------------------------------------
%
\begin{frame}{String rewriting}
  Let us consider a string of white and black beads, like \(\circ
  \bullet \bullet \bullet \circ \circ \bullet\), and a one-player game
  whose unique move consists in replacing two adjacent beads with only
  one according to the rules
  \begin{equation*}
    \bullet \; \circ   \xrightarrow{\alpha} \bullet\qquad\qquad
    \circ   \; \bullet \xrightarrow{\beta} \bullet\qquad\qquad
    \bullet \; \bullet \xrightarrow{\gamma} \circ
  \end{equation*}

  The rules~\(\alpha\), \(\beta\)~and~\(\gamma\) make up a simple
  \textbf{string\hyp{}rewriting system}.

  \bigskip

  Rules \(\alpha\)~and~\(\beta\) mean: \emph{`A black bead absorbs the
  white bead next to it.'}

  \bigskip

  The goal of this game is to end up with as few beads as possible, so
  our example may lead to the \textbf{rewrites}
  \begin{equation*}
    \circ \bullet \bullet \, \fbox{\(\bullet \, \circ\)} \circ \bullet
    \xrightarrow{\alpha} \circ \bullet \bullet \, \fbox{\(\bullet \,
      \circ\)} \, \bullet \xrightarrow{\alpha} \fbox{\(\circ \,
      \bullet\)} \bullet \bullet \bullet \xrightarrow{\beta} \bullet
    \bullet \fbox{\(\bullet \, \bullet\)} \xrightarrow{\gamma} \bullet \,
    \fbox{\(\bullet \, \circ\)} \xrightarrow{\alpha} \fbox{\(\bullet
      \, \bullet\)} \xrightarrow{\gamma} \circ
  \end{equation*}

  \medskip

  \noindent where the part of the string to be rewritten next is
  framed.

\end{frame}

% ------------------------------------------------------------------------
%
\begin{frame}{Normal forms and termination}

  \label{termination1}

  Other compositions of the rules lead to the same result~\(\circ\) as
  well. Some others bring all\hyp{}white or all\hyp{}black strings,
  like \(\circ \, \circ\) and~\(\bullet\).

  \bigskip

  Strings that can not be further rewritten, or \textbf{reduced}, are
  called \textbf{normal forms}. We may wonder whether all strings have
  a normal form, and, if so, if it is unique and, furthermore, if it
  is either all\hyp{}white or all\hyp{}black.

  \bigskip

  First, let us note that the system is \textbf{terminating}, that is,
  there is no infinite chain of rewrites, because the number of beads
  strictly decreases in all the rules.

  \bigskip

  This is not a necessary condition in general, for instance, \(\circ
  \; \bullet \xrightarrow{\beta} \bullet \circ \circ\) would preserve
  termination because the composition \(\beta\alpha\alpha\) would be
  equivalent to the original rule~\(\beta\).

  \bigskip

  In particular, this means that any string has a normal form.

\end{frame}

% ------------------------------------------------------------------------
%
\begin{frame}{Invariants}

  Second, notice how the parity of the number of black beads is
  \textbf{invariant} through each rule and how there is no rewrite
  rule for two adjacent white beads.

  \bigskip

  Therefore, if there are \(2p\) initial black beads, then composing
  rules \(\alpha\)~and~\(\beta\) lead to an all\hyp{}black string,
  which can be reduced by applying rule~\(\gamma\) to contiguous pairs
  of beads into an all\hyp{}white string made of \(p\)~beads.

  \bigskip

  Otherwise, the same all\hyp{}black string can be reduced by applying
  alternatively \(\gamma\)~and~\(\beta\) on the left end or
  \(\gamma\)~and~\(\alpha\) on the right end, yielding~\(\circ\).

  \bigskip

  Similarly, if there is an initial odd number of black beads, we
  always end up with one black bead.

\end{frame}

% ------------------------------------------------------------------------
%
\begin{frame}{Confluence}

  Consequently, normal forms are not unique. A simple example is
  \begin{equation*}
    \circ \circ \xleftarrow{\gamma} \bullet \bullet \circ
    \xrightarrow{\alpha} \bullet \bullet \xrightarrow{\gamma}
    \circ
  \end{equation*}

  \medskip

  Systems where normal forms are unique are \textbf{confluent}.

  \bigskip

  If we think of a rewrite step as an elementary computation, it
  becomes clear that confluence is a desirable property for carrying
  out computations, as we would expect, for instance, an arithmetic
  expression to have the same value no matter the order of the small
  steps.

  \bigskip

  A systemic order of rule application is called a \textbf{reduction
    strategy}.

\end{frame}

% ------------------------------------------------------------------------
%
\begin{frame}{Completing a system}

  We can obtain confluence by adding the rule~\(\delta\):
  \begin{equation*}
    \bullet \; \circ   \xrightarrow{\alpha} \bullet\qquad\qquad
    \circ   \; \bullet \xrightarrow{\beta} \bullet\qquad\qquad
    \bullet \; \bullet \xrightarrow{\gamma}
    \circ\qquad\qquad
    \circ\; \circ \xrightarrow{\delta} \circ
  \end{equation*}
  The result of the game is now always one bead whose colour depends
  on the original parity of the black beads as before.

  \bigskip

  To see why, let us consider first that two non\hyp{}overlapping
  parts of a string can be rewritten in parallel, so they are not a
  concern.

  \bigskip

  The interesting cases occur when two applications of rules (maybe of
  the same rule) lead to different strings because they do
  overlap. For instance,
  \begin{equation*}
    \circ \, \circ \xleftarrow{\gamma} \bullet \bullet
    \circ \xrightarrow{\alpha} \bullet \, \bullet
  \end{equation*}

  \medskip

  The important point is that \(\circ \, \circ\) and \(\bullet \,
  \bullet\) can be rewritten into~\(\circ\) at the next step by
  \(\delta\)~and~\(\gamma\), respectively.

\end{frame}

% ------------------------------------------------------------------------
%
\begin{frame}{Critical pairs}

  In general, what matters is that all pairs of strings resulting from
  the application of overlapping rules, called \textbf{critical
    pairs}, can be rewritten to the same string: they are
  \textbf{joinable}.

  \bigskip

  In our example, all interactions occur on substrings made of three
  beads, so we must examine eight cases:
  \begin{center}
    \includegraphics[]{bullets}
  \end{center}

  In all the cases, the divergences are joinable in one step at most.

\end{frame}

% ------------------------------------------------------------------------
%
\begin{frame}{Local confluence}

  If all critical pairs are joinable after a divergence in one ore
  more rewrites, the system is \textbf{locally confluent}.

  \bigskip

  Moreover, if the system is also terminating, then \textbf{every
    string has exactly one normal form}.

  \bigskip

  Sometimes, as seen previously, a non\hyp{}confluent system can be
  completed to become confluent, via a \textbf{completion procedure},
  for example, in the case of terminating systems, the
  \textbf{Knuth\hyp{}Bendix semi\hyp{}algorithm}.

\end{frame}

% ------------------------------------------------------------------------
%
\begin{frame}{Ground and linear rewrite systems}

  The system we defined is \textbf{ground}, that is, it involves no
  variables.

  \bigskip

  These allow a finite system to denote an infinite number of ground
  rules if variables range over infinite sets, or simply, they reduce
  the size of rewrite systems.

  \bigskip

  For instance, our continued example is equivalent to
  \begin{equation*}
    \bullet \; \circ \xrightarrow{\alpha} \bullet
    \qquad\qquad
    \circ \; x \xrightarrow{\beta+\delta} x
    \qquad\qquad
    \bullet \; \bullet \xrightarrow{\gamma} \circ
  \end{equation*}

  \medskip

  If we accept multiple occurrences of a variable on the
  left\hyp{}hand side of a rule (that is, not left\hyp{}linear), we
  can further decrease the size of the system:
  \begin{equation*}
    x  \; x \xrightarrow{\gamma+\delta} \circ\qquad\qquad
    x  \; y \xrightarrow{\alpha+\beta} \bullet
  \end{equation*}

\end{frame}

% ------------------------------------------------------------------------
%
\begin{frame}{Ordered rewrite systems}

  Notice that there is now an implicit order over the rules: the rule
  \(\gamma+\delta\) must be examined first for a \textbf{match} with a
  part of the current string, because it is included in the second
  (set \(x=y\) in \(\alpha+\beta\)).

  \bigskip

  Non\hyp{}linear systems are equational theories due the implicit
  equality on substructures, which is complex: we will limit ourselves
  to \textbf{linear systems}.

  \bigskip

  In general, the order in which the rules are written on the page is
  their \textbf{logical order}.

  \bigskip

  Moreover, let us remark that the system does not specify that
  \(x\)~must either be~\((\circ)\) or~\((\bullet)\): this should be
  stated nevertheless somewhere else.

  \bigskip

  Generally speaking, this means that the \textbf{type} of the
  variables has to be defined elsewhere or inferred from the variable
  uses.

\end{frame}
